\documentclass{article}
\usepackage[brazilian]{babel}
\usepackage[T1]{fontenc}
\usepackage[utf8]{inputenc}
\usepackage{amsmath}

\title{Introdu\c{c}\~ao \`a teoria dos grafos}
\author{P.~Silva}
\date{2026}

\begin{document}
\maketitle

\section{Introdu\c{c}\~ao}

A teoria dos grafos \'e um ramo da matem\'atica que estuda as rela\c{c}\~oes
entre objetos.  Um grafo $G = (V, E)$ consiste em um conjunto de v\'ertices~$V$
e um conjunto de arestas~$E$.

Neste trabalho, estudamos a colora\c{c}\~ao de grafos planares.  O teorema das
quatro cores afirma que todo grafo planar pode ser colorido com no m\'aximo
quatro cores de modo que v\'ertices adjacentes recebam cores distintas.

\section{Resultados}

\begin{theorem}
Para todo grafo planar $G$, temos $\chi(G) \le 4$, onde $\chi(G)$ denota o
n\'umero crom\'atico de~$G$.
\end{theorem}

A demonstra\c{c}\~ao original de Appel e Haken utiliza verifica\c{c}\~ao
computacional de 1.936~configura\c{c}\~oes redut\'iveis.  Uma prova
simplificada foi obtida por Robertson, Sanders, Seymour e Thomas em 1997.

\section{Conclus\~ao}

Os m\'etodos computacionais s\~ao essenciais na teoria dos grafos moderna.

\end{document}
